% Independent local nontraceability proof; not part of the main manuscript.
Let $B$ be the listed pentagon graph with boundary $00001$.
Orient the three colours cyclically $1\to2\to3\to1$, write
$\epsilon(a,b)=1$ in this direction and $-1$ in the other, and put
\[
 \sigma(w)=\epsilon(w_1,w_2)\epsilon(w_3,w_4),\qquad
 \chi(w)=\begin{cases}\sigma(w),&w_5=0,\\-\sigma(w),&w_5\in\{2,3\}.
 \end{cases}
\]
Every move at the first four positions flips $\chi$, as does a last-position
move involving $0$. A last-position move between $2$ and $3$ requires
$w_1=w_4=1$, so the only same-sign edges are
$12312\!-12313$ and $13212\!-13213$, both negative.
The $24$-state core $w_5=0$ has twelve states of each sign, paired by changing
$w_1$ to the other colour different from $w_2$.
The remaining states form the following two paths:
\[
\begin{array}{l}
23123,13123,12123,12323,12313,12312,32312,31312,31212,31232,\\
21323,21313,21213,23213,13213,13212,13232,13132,12132,32132.
\end{array}
\]
Their central edges are the two negative same-sign edges, and every other
edge flips sign. Each path therefore has six negative and four positive
states. Thus $|\chi^{-1}(-1)|=24$ and $|\chi^{-1}(1)|=20$.
Deleting $S=\chi^{-1}(1)$ leaves two disjoint edges and twenty isolated
vertices, so $c(B-S)=22>|S|+1=21$. Since deleting $S$ from a spanning path
leaves at most $|S|+1$ components, $B$ has no Hamilton path.
